\documentclass[../../main.tex]{subfiles}
\begin{document}

\begin{center}
  {\fontsize{48pt}{58pt}\selectfont \textbf{1991}\normalsize\textbf{年度}}
\end{center}

\begin{table}[htb]
  \centering
  \caption{出題テーマと難易度}
  \begin{tabular}{|c|p{9cm}|c|}
    \hline
    \textbf{問題番号} & \textbf{テーマ} & \textbf{難易度} \\
    \hline
    第1問 & 円の交点から定まる放物線の軌跡と接線 & \\
    \hline
    第2問 & 行列による一次変換と射影 & \\
    \hline
    第3問 & 対辺が垂直な四面体の中点と球面 & \\
    \hline
    第4問 & 正接関数の相加平均・相乗平均不等式 & \\
    \hline
    第5問 & 玉を取り出して補充する試行の確率 & \\
    \hline
    第6問 & 積分方程式から定まる関数の導関数 & \\
    \hline
  \end{tabular}
\end{table}

\begin{table}[htb]
  \centering
  \caption{解答}
  \renewcommand{\arraystretch}{1.4}
  \begin{tabular}{|c|c|p{9cm}|}
    \hline
    \textbf{問題番号} & \textbf{小問} & \textbf{答え} \\
    \hline
    第1問 & (1) & $s=2\sqrt{p+1}$，$t=p+2$ \\
    \cline{2-3}
    & (2) & $y=\dfrac{x^2}{4}+1$（$x>2$），$PQ$は接線（証明問題） \\
    \hline
    第2問 & (1) &  \\
    \cline{2-3}
    & (2) &  \\
    \cline{2-3}
    & (3) &  \\
    \hline
    第3問 & -- & 6本の辺の中点は重心を中心とする同一球面上にある（証明問題） \\
    \hline
    第4問 & -- & $\sqrt{\tan a\tan b}\le\tan\left(\dfrac{a+b}{2}\right)\le\dfrac{\tan a+\tan b}{2}$（証明問題） \\
    \hline
    第5問 & (1) & $P_{1,n+1}=\dfrac{(N+2)P_{1,n}+P_{2,n}}{N+3}$，$P_{2,n+1}=\dfrac{(N+1)P_{2,n}+2P_{3,n}}{N+3}$，$P_{3,n+1}=\dfrac{N P_{3,n}}{N+3}$ \\
    \cline{2-3}
    & (2) & $P_{n+1}=\dfrac{N+2}{N+3}P_n$，$P_n=\dfrac{3}{N+3}\left(\dfrac{N+2}{N+3}\right)^{n-1}$ \\
    \hline
    第6問 & (1) & $f(0)=a$，$f'(x)=\dfrac{1}{1+f(x)^a e^{-f(x)^2/2}}$ \\
    \cline{2-3}
    & (2) & $0<b\le f'(x)\le1$，$0\le f(x)(1/f'(x)-1)\le c$となる$b,c$が存在 \\
    \cline{2-3}
    & (3) & $\lim_{x\to\infty}f'(x)=1$；最小値は$a\le1$で$\dfrac{e^{a/2}}{a^{a/2}+e^{a/2}}$，$a\ge1$で$\dfrac{e^{a^2/2}}{a^a+e^{a^2/2}}$ \\
    \hline
  \end{tabular}
\end{table}

\paragraph{試験形態}
6問．試験時間は要確認．

\end{document}
